\documentclass[10pt,a4paper]{amsart}
\usepackage[margin=19mm]{geometry}
\usepackage{amsmath,amssymb,mathtools}
\usepackage[scr=rsfs,cal=euler]{mathalfa}
\usepackage{xfrac}
\usepackage[unicode,hidelinks,bookmarks=false]{hyperref}
\newcommand{\ie}{i.e.\ }
\newcommand{\cf}{cf.\ }
\newcommand{\resp}{respectively\ }
\newcommand{\Subsection}{\subsection}
\providecommand{\nobreakdash}{\nobreak\hbox{-}}
\setlength{\emergencystretch}{2em}
\multlinegap0pt
\usepackage{bm}
\usepackage{bbm}
\usepackage{dsfont}
\usepackage{xspace}
\usepackage{cancel}
\usepackage{mathtools}
\usepackage{amsthm}
\makeatletter
\@ifundefined{lemma}{\newtheorem{lemma}{Lemma}}{}
\@ifundefined{proposition}{\newtheorem{proposition}{Proposition}}{}
\makeatother
\newcommand{\HS}{\mathrm{HS}}
\newcommand{\Id}{\mathrm{Id}}
\newcommand{\KL}{\mathrm{KL}}
\newcommand{\Law}{\mathcal{L}}
\newcommand{\Prob}{\mathbb{P}}
\newcommand{\R}{\mathbb{R}}
\newcommand{\SO}{\mathrm{SO}}
\newcommand{\TV}{\mathrm{TV}}
\newcommand{\dd}{\mathrm{d}}
\newcommand{\norm}[1]{\left\lVert #1\right\rVert}
\newcommand{\op}{\mathrm{op}}
\newcommand{\so}{\mathfrak{so}}
\title[Kac's walk on rotation matrices mixes in $n^2 \log n$ steps]{Kac's walk on rotation matrices mixes in $n^2 \log n$ steps}
\author{Natesh S. Pillai, Aaron Smith}
\date{Source: arXiv:2604.23828v1 (CC BY 4.0)}
\begin{document}
\maketitle

\noindent\emph{Source and licence.} Kac's walk on rotation matrices mixes in $n^2 \log n$ steps. Version arXiv:2604.23828v1 (CC BY 4.0), distributed under the Creative Commons Attribution 4.0 International licence, as recorded in the arXiv licence metadata for this exact version and shown on the abstract page https://arxiv.org/abs/2604.23828v1. Full rights notice: licenses/arxiv-2604.23828-CC-BY-4.0.txt.\par\smallskip

\noindent\textit{Editorial scope.} Counted unit: the Proposition whose source label is \texttt{prop:gaussian-bch} in the article (source0.tex lines 1734--1768, source label \texttt{prop:gaussian-bch}), delivered together with the article's own complete original proof of it (source0.tex lines 1769--2139, one \texttt{proof} environment of 371 lines). No mathematics is written, repaired or re-derived: statement and proof are the article's own text line for line. Required context retained uncounted: prop:near-identity (lines 1162--1186); lem:radial-cutoff (lines 1295--1311); lem:gaussian-shift (lines 1625--1635); lem:bch-remainder (lines 1653--1675). Every definition, display and label of those lines is retained unchanged. Omitted from this package and not redistributed: 1--1161, 1187--1294, 1312--1624, 1636--1652, 1676--1733, 2140--2651. Nothing inside a retained excerpt is omitted or altered: every retained body is delivered exactly as the article's lines are written, so no omission receipt of any kind accompanies this package. Citations: the retained excerpts carry no citation command at all, so no bibliography entry of the article is transcribed and no reference is redistributed. Reference adaptation: none was needed and none was made; every reference inside the delivered unit resolves inside it, so no unresolved reference or citation is produced. Printed slips kept literal: no printed word, symbol, hypothesis or number of the counted statement or of its proof was altered. Layout and class: the article's own class is \texttt{amsart} with packages including pstricks, times, comment, xcolor, natbib, hyperref, amsmath, amssymb, amsthm, mathtools, enumitem, inputenc, amsfonts, bbm; the wrapper is the lane's pinned \texttt{amsart} editorial wrapper at 10pt on A4, which restates the theorem-like environments the retained text needs and adds the article's own macro definitions that the retained text needs (\texttt{\textbackslash HS}, \texttt{\textbackslash Id}, \texttt{\textbackslash KL}, \texttt{\textbackslash Law}, \texttt{\textbackslash Prob}, \texttt{\textbackslash R}, \texttt{\textbackslash SO}, \texttt{\textbackslash TV}, \texttt{\textbackslash dd}, \texttt{\textbackslash norm}, \texttt{\textbackslash op}, \texttt{\textbackslash so}), with extra wrapper packages (bm, bbm, dsfont, xspace, cancel, mathtools). The delivered numbering is the unit's own. Assets: no retained span contains a figure environment, a graphics-inclusion command, a PDF-page inclusion, a table or a TikZ picture; no image file is copied into this package, the package declares no asset dependency and no image file is redistributed with it. Licence: version arXiv:2604.23828v1 of this article is distributed under the Creative Commons Attribution 4.0 International licence, as recorded in the arXiv licence metadata for this exact version and shown on the abstract page https://arxiv.org/abs/2604.23828v1. A full rights notice is delivered at licenses/arxiv-2604.23828-CC-BY-4.0.txt. Characters: every retained excerpt is pure ASCII, so the delivered engine, XeLaTeX with its default Latin Modern text font, reports no missing character.
\begin{proposition}[Gaussian stability under small perturbations]
\label{prop:near-identity}
Let $\gamma_d=\mathcal N(0,I_d)$. Let $R:\R^d\to\R^d$ be a $C^1$ map satisfying
\[
\sup_{x\in\R^d}\|R(x)\|_2\le \alpha,
\qquad
\sup_{x\in\R^d}\|\dd R(x)\|_{\op}\le \beta
\]
with $0<\beta\le 1/2$. Set $T=\Id+R$. Then $T$ is a global $C^1$
diffeomorphism, and there exists a universal constant $C_{\mathrm{NI}}>0$ such
that
\begin{equation}
\label{eq:near-identity-kl}
\KL(T_\#\gamma_d\,\|\,\gamma_d)
\le
C_{\mathrm{NI}}\bigl(\alpha^2+d\beta^2\bigr),
\end{equation}
and
\begin{equation}
\label{eq:near-identity-tv}
\|T_\#\gamma_d-\gamma_d\|_{\TV}
\le
C_{\mathrm{NI}}\bigl(\alpha+\sqrt d\,\beta\bigr).
\end{equation}
\end{proposition}

\begin{lemma}[Radial cutoff]
\label{lem:radial-cutoff}
There is an absolute constant $C_{\mathrm{cut}}<\infty$ such that for every
$d\ge 1$ and every $R>0$, there exists a smooth function
$\Psi_{d,R}:\R^d\to[0,1]$ satisfying
\[
\Psi_{d,R}(x)=1 \quad \text{if } \|x\|_2\le 2R,
\qquad
\Psi_{d,R}(x)=0 \quad \text{if } \|x\|_2\ge 3R,
\]
and
\[
\sup_{x\in\R^d}\|\nabla\Psi_{d,R}(x)\|_2
\le
\frac{C_{\mathrm{cut}}}{R}.
\]
\end{lemma}

\begin{lemma}
\label{lem:gaussian-shift}
If $G\sim \mathcal N(0,\Sigma)$ on $\R^N$ with $\Sigma$ positive definite, then for every
$h\in\R^N$,
\begin{equation}
\label{eq:gaussian-shift}
\norm{\Law(G+h)-\Law(G)}_{\TV}
\le
\frac{1}{\sqrt{2}}\norm{\Sigma^{-1/2}h}_2.
\end{equation}
\end{lemma}

\begin{lemma}
\label{lem:bch-remainder}
There exist constants $r_\ast,C_\ast>0$ independent of $n$ such that the map
\[
\Phi(\xi,h) \equiv \log\!\bigl(\exp(\xi)\exp(h)\bigr)
\]
is smooth on
\[
\{(\xi,h)\in \so(n)\times \so(n): \norm{\xi}_{\HS}\le r_\ast,\ \norm{h}_{\HS}\le r_\ast\},
\]
and the remainder
\begin{equation} \label{EqDefRemainderLemma6}
\beta_h(\xi)=\Phi(\xi,h)-\xi-h
\end{equation}
satisfies
\begin{equation}
\label{eq:bch-remainder}
\norm{\beta_h(\xi)}_{\HS}\le C_\ast \norm{\xi}_{\HS}\norm{h}_{\HS},
\qquad
\norm{\dd_\xi \beta_h(\xi)}_{\op}\le C_\ast \norm{h}_{\HS}
\end{equation}
whenever $\norm{\xi}_{\HS},\norm{h}_{\HS}\le r_\ast$.
\end{lemma}

\begin{proposition}[Gaussian overlap transfer on $\SO(n)$]
\label{prop:gaussian-bch}
Let $\mu$ be a probability measure on $\so(n)$ supported on
\begin{equation} \label{SupportHyp1}
\{\xi\in \so(n): \norm{\xi}_{\HS}\le r_\ast/4\},
\end{equation}
 where $r_\ast$ is the radius from Lemma~\ref{lem:bch-remainder},
and let $\nu=\exp_\#\mu.$
Let $\gamma=\mathcal N(0,\Sigma)$
be a centered Gaussian law on $\so(n)$, where $\Sigma$ is positive definite, and
assume
\begin{equation}
\label{eq:gaussian-transfer-chart}
\norm{\Sigma^{1/2}}_{\op}\le \frac{r_\ast}{16\sqrt N}.
\end{equation}
Set $\kappa(\Sigma)=\frac{\lambda_{\max}(\Sigma)}{\lambda_{\min}(\Sigma)}$.
There exist absolute constants $c_{\mathrm{GT}},c_B,C_B>0$ with the following
property: for every $h\in \so(n)$ satisfying
\begin{equation}
\label{eq:gaussian-transfer-h-small}
\norm{h}_{\HS}\le c_{\mathrm{GT}}\,\kappa(\Sigma)^{-1/2},
\end{equation}
and with $g=\exp(h)$, 
\begin{align}
\norm{\nu-(T_g)_\#\nu}_{\TV}
\le
&2\norm{\mu-\gamma}_{\TV}
+
C_B\norm{\Sigma^{-1/2}h}_{\HS} \notag \\
&+
C_B\sqrt N\,\kappa(\Sigma)^{1/2}\norm{h}_{\HS}
+
e^{-c_BN}.\label{eq:gaussian-bch}
\end{align}
\end{proposition}

\begin{proof}
Let \(r_\ast\) and \(C_\ast\) be the constants from
Lemma~\ref{lem:bch-remainder}. Define
\[
c_{\mathrm{GT}}
:=
\min\left\{
r_\ast,\,
\frac{1}{2C_\ast(1+3C_{\mathrm{cut}})}
\right\}.
\]
Since \(\kappa(\Sigma)\ge 1\), the assumption
\[
\|h\|_{\HS}\le c_{\mathrm{GT}}\kappa(\Sigma)^{-1/2}
\]
implies
\[
\|h\|_{\HS}\le c_{\mathrm{GT}}\le r_\ast.
\]

Choose a measurable extension
\[
B_h:\so(n)\to\so(n)
\]
of the local map
\[
\xi\mapsto \log(\exp(\xi)\exp(h))
\]
from the ball \(\{\|\xi\|_{\HS}\le r_\ast/4\}\) to all of \(\so(n)\).
Since \(\mu\) is supported in this ball, we have, for \(\xi\in \mathrm{supp}\,\mu\),
\[
B_h(\xi)=\log(\exp(\xi)\exp(h)).
\]
Thus
\[
T_g(\exp(\xi))
=
\exp(\xi)\exp(h)
=
\exp(B_h(\xi)),
\]
and hence
\[
\nu=\exp_\#\mu,
\qquad
(T_g)_\#\nu=\exp_\#(B_h)_\#\mu.
\]
Since total variation cannot increase under pushforward,
\[
\|\nu-(T_g)_\#\nu\|_{\TV}
\le
\|\mu-(B_h)_\#\mu\|_{\TV}.
\]
Inserting \(\gamma\) and \((B_h)_\#\gamma\), and then using contraction of
total variation under the pushforward by the measurable map \(B_h\), gives
\begin{align}
\|\mu-(B_h)_\#\mu\|_{\TV}
&\le
\|\mu-\gamma\|_{\TV}
+
\|\gamma-(B_h)_\#\gamma\|_{\TV}
+
\|(B_h)_\#\gamma-(B_h)_\#\mu\|_{\TV}
\notag\\
&\le
\|\mu-\gamma\|_{\TV}
+
\|\gamma-(B_h)_\#\gamma\|_{\TV}
+
\|\gamma-\mu\|_{\TV}
\notag\\
&=
2\|\mu-\gamma\|_{\TV}
+
\|\gamma-(B_h)_\#\gamma\|_{\TV}.
\label{eq:transfer-triangle}
\end{align}

It remains to control the second term. We split it into a mean shift and a nonlinear Baker-Campbell-Hausdorff correction; the latter is localized and whitened so that Proposition~\ref{prop:near-identity} can be applied globally. Let
\[
\tau_h(\xi)=\xi+h.
\]
We have
\begin{equation}
\label{eq:transfer-split}
\|\gamma-(B_h)_\#\gamma\|_{\TV}
\le
\|\gamma-(\tau_h)_\#\gamma\|_{\TV}
+
\|(\tau_h)_\#\gamma-(B_h)_\#\gamma\|_{\TV}.
\end{equation}
By Lemma~\ref{lem:gaussian-shift},
\begin{equation}
\label{eq:transfer-gaussian-shift}
\|\gamma-(\tau_h)_\#\gamma\|_{\TV}
\le
\frac1{\sqrt2}\|\Sigma^{-1/2}h\|_{\HS}.
\end{equation}
This controls the first term in \eqref{eq:transfer-split}. We now bound the nonlinear term. Set
\[
\kappa:=\kappa(\Sigma)
=
\frac{\lambda_{\max}(\Sigma)}{\lambda_{\min}(\Sigma)}.
\]
Write
\[
\Xi=\Sigma^{1/2}X,
\qquad
X\sim\mathcal N(0,I_N),
\]
so that \(\Xi\sim\gamma\). Let
\[
F=\{\|X\|_2\le 2\sqrt N\}.
\]
A standard Chernoff bound gives
\begin{equation}
\label{eq:transfer-X-tail}
\Prob(F^c)\le e^{-2c_BN},
\end{equation}
where
\[
c_B:=\frac{3-\log 4}{4}.
\]
Using the construction and notation from the statement of Lemma~\ref{lem:radial-cutoff}, choose
\[
\Psi=\Psi_{N,\sqrt N}.
\]
Then
\[
\Psi(x)=1 \quad\text{if }\|x\|_2\le 2\sqrt N,
\qquad
\Psi(x)=0 \quad\text{if }\|x\|_2\ge 3\sqrt N,
\]
and
\begin{equation}
\label{eq:transfer-cutoff-grad}
\sup_{x\in\R^N}\|\nabla\Psi(x)\|_2
\le
C_{\mathrm{cut}}N^{-1/2}.
\end{equation}
Recall that $r_\ast$ is the radius from Lemma~\ref{lem:bch-remainder}. We first check that the whitened variable stays in this logarithmic chart on the ball that contains the support of the cutoff.
By the chart assumption \eqref{eq:gaussian-transfer-chart}, if
\(\|x\|_2\le 4\sqrt N\), then
\[
\|\Sigma^{1/2}x\|_{\HS}
\le
\|\Sigma^{1/2}\|_{\op}\|x\|_2
\le
\frac{r_\ast}{4}.
\]
Thus, on this ball, the map
\[
x\mapsto \log\!\bigl(\exp(\Sigma^{1/2}x)\exp(h)\bigr)
\]
is well-defined and smooth.

We next write the Baker-Campbell-Hausdorff remainder in whitened coordinates and then use the cutoff function to obtain the global bounds required by Proposition~\ref{prop:near-identity}.
Define the local whitened error
\[
b_h(x):=\Sigma^{-1/2}\beta_h(\Sigma^{1/2}x),
\qquad \|x\|_2\le 4\sqrt N,
\]
where
\[
\beta_h(\xi)=\log(\exp(\xi)\exp(h))-\xi-h
\]
is the local remainder from Lemma~\ref{lem:bch-remainder}. 
By the chain rule,
\[
\dd b_h(x)
=
\Sigma^{-1/2}\circ
\dd_\xi\beta_h(\Sigma^{1/2}x)
\circ \Sigma^{1/2}.
\]
Since
\[
\|\Sigma^{-1/2}\|_{\op}\|\Sigma^{1/2}\|_{\op}
=
\sqrt{\kappa},
\]
Lemma~\ref{lem:bch-remainder} gives, for \(\|x\|_2\le 4\sqrt N\),
\begin{equation}
\label{eq:transfer-b-bound}
\|b_h(x)\|_2
\le
C_\ast\sqrt{\kappa}\,\|x\|_2\|h\|_{\HS},
\end{equation}
and
\begin{equation}
\label{eq:transfer-db-bound}
\|\dd b_h(x)\|_{\op}
\le
C_\ast\sqrt{\kappa}\,\|h\|_{\HS}.
\end{equation}

Define the global cutoff error
\[
R_h(x):=
\begin{cases}
\Psi(x)b_h(x), & \|x\|_2\le 4\sqrt N,\\
0, & \|x\|_2>4\sqrt N,
\end{cases}
\qquad
S_h(x):=x+R_h(x).
\]
$S_{h}$ is well-defined and \(C^1\), because \(\Psi(x)=0\) when
\(\|x\|_2\ge 3\sqrt N\). Since \(0\le \Psi\le 1\) and \(\Psi\) is supported in
\(\{\|x\|_2\le 3\sqrt N\}\), \eqref{eq:transfer-b-bound} gives
\begin{equation}
\label{eq:transfer-R-bound}
\sup_{x\in\R^N}\|R_h(x)\|_2
\le
3C_\ast\sqrt{\kappa N}\,\|h\|_{\HS}.
\end{equation}

We next bound the derivative. On the support of \(\Psi\),
\[
\dd R_h(x)
=
\Psi(x)\dd b_h(x)+(\nabla\Psi(x))\otimes b_h(x),
\]
and \(\dd R_h(x)=0\) outside this support. The first term is controlled by
\eqref{eq:transfer-db-bound}:
\[
\|\Psi(x)\dd b_h(x)\|_{\op}
\le
C_\ast\sqrt{\kappa}\,\|h\|_{\HS}.
\]
For the second term, on the support of \(\nabla\Psi\) we have
\(\|x\|_2\le 3\sqrt N\). Therefore \eqref{eq:transfer-cutoff-grad} and
\eqref{eq:transfer-b-bound} imply
\[
\|(\nabla\Psi(x))\otimes b_h(x)\|_{\op}
\le
C_{\mathrm{cut}}N^{-1/2}
\cdot
3C_\ast\sqrt{\kappa N}\,\|h\|_{\HS}.
\]
Combining the two estimates gives
\begin{equation}
\label{eq:transfer-DR-bound}
\sup_{x\in\R^N}\|\dd R_h(x)\|_{\op}
\le
C_\ast(1+3C_{\mathrm{cut}})\sqrt{\kappa}\,\|h\|_{\HS}.
\end{equation}

By the definition of \(c_{\mathrm{GT}}\), the assumption on \(h\) implies
\[
\sup_{x\in\R^N}\|\dd R_h(x)\|_{\op}\le \frac12.
\]
Therefore Proposition~\ref{prop:near-identity} applies to
\[
S_h=\Id+R_h.
\]
With
\[
\alpha=3C_\ast\sqrt{\kappa N}\,\|h\|_{\HS},
\qquad
\beta=C_\ast(1+3C_{\mathrm{cut}})\sqrt{\kappa}\,\|h\|_{\HS},
\]
it yields
\begin{equation}
\label{eq:transfer-Sh-close}
\|(S_h)_\#\mathcal N(0,I_N)-\mathcal N(0,I_N)\|_{\TV}
\le
C_{\mathrm{NI}}C_\ast(4+3C_{\mathrm{cut}})
\sqrt{\kappa N}\,\|h\|_{\HS}.
\end{equation}


We have now controlled the whitened cutoff map \(S_h=\Id+R_h\). It remains to
compare this cutoff model with the original map \(B_h\). Define
\[
\widetilde B_h(x)=\Sigma^{-1/2}B_h(\Sigma^{1/2}x),
\qquad
u=\Sigma^{-1/2}h.
\]
On the event \(F\), the cutoff is inactive: \(\Psi(X)=1\). Moreover,
\(\|\Sigma^{1/2}X\|_{\HS}\le r_\ast/4\), so \(B_h\) agrees with the local
logarithmic map at \(\Sigma^{1/2}X\). Hence, on \(F\),
\[
B_h(\Sigma^{1/2}X)
=
\log\!\bigl(\exp(\Sigma^{1/2}X)\exp(h)\bigr)
=
\Sigma^{1/2}X+h+\beta_h(\Sigma^{1/2}X),
\]
and therefore
\[
\widetilde B_h(X)
=
X+u+b_h(X)
=
\tau_u(S_h(X)),
\qquad
\tau_u(x)=x+u.
\]
Thus the coupling inequality gives
\[
\|\Law(\widetilde B_h(X))-\Law(\tau_u(S_h(X)))\|_{\TV}
\le
\Prob(F^c).
\]
By the triangle inequality and translation invariance of total variation,
\begin{align*}
\|\Law(\widetilde B_h(X))-\Law(\tau_u(X))\|_{\TV}
&\le
\Prob(F^c)
+
\|\Law(\tau_u(S_h(X)))-\Law(\tau_u(X))\|_{\TV}\\
&=
\Prob(F^c)
+
\|(S_h)_\#\mathcal N(0,I_N)-\mathcal N(0,I_N)\|_{\TV}.
\end{align*}
Using \eqref{eq:transfer-X-tail} in the form
\(\Prob(F^c)\le e^{-2c_BN}\), together with
\eqref{eq:transfer-Sh-close}, gives
\begin{align}
\|(\widetilde B_h)_\#\mathcal N(0,I_N)
-(\tau_u)_\#\mathcal N(0,I_N)\|_{\TV}
&\le
C_{\mathrm{NI}}C_\ast(4+3C_{\mathrm{cut}})
\sqrt{\kappa N}\,\|h\|_{\HS}
+
e^{-2c_BN}.
\label{eq:transfer-whitened-nonlinear}
\end{align}
Finally, applying the invertible linear map \(\Sigma^{1/2}\) to both measures and
using
\[
B_h\circ \Sigma^{1/2}=\Sigma^{1/2}\circ \widetilde B_h,
\qquad
\tau_h\circ \Sigma^{1/2}=\Sigma^{1/2}\circ \tau_u,
\]
turns \eqref{eq:transfer-whitened-nonlinear} into
\begin{equation}
\label{eq:transfer-nonlinear}
\|(B_h)_\#\gamma-(\tau_h)_\#\gamma\|_{\TV}
\le
C_{\mathrm{NI}}C_\ast(4+3C_{\mathrm{cut}})
\sqrt{\kappa N}\,\|h\|_{\HS}
+
e^{-2c_BN}.
\end{equation}


Finally set
\[
C_B:=
\max\left\{
\frac1{\sqrt2},
C_{\mathrm{NI}}C_\ast(4+3C_{\mathrm{cut}})
\right\}.
\]
Combining \eqref{eq:transfer-triangle}, \eqref{eq:transfer-split},
\eqref{eq:transfer-gaussian-shift}, and \eqref{eq:transfer-nonlinear}, we get
\[
\|\nu-(T_g)_\#\nu\|_{\TV}
\le
2\|\mu-\gamma\|_{\TV}
+
C_B\|\Sigma^{-1/2}h\|_{\HS}
+
C_B\sqrt{\kappa N}\,\|h\|_{\HS}
+
e^{-c_BN}.
\]
Since \(\kappa=\kappa(\Sigma)\), this is exactly \eqref{eq:gaussian-bch} and the proof is finished.
\end{proof}

\end{document}
